\section{Four-state receipt recovery}

The blind census also contains a unique normal \(\Vfour\) class.  It is
class \(20\), generated by registered permutations \(65\) and \(124\).
Their product is permutation \(326\), the already recovered normal
binary deck transformation.  Thus the blind actions themselves give
the subgroup chain
\[
  \langle326\rangle
  \cong\Ctwo
  \leq
  \langle65,124\rangle
  \cong\Vfour.
\]

The normal \(\Vfour\) has fifteen vertex orbits, each of size four.
The orbit quotient has fifteen vertices, thirty edges, degree four,
no loops, and four native lifts above every quotient edge.  Its labeled
edge set agrees exactly with the independently preserved
\(\Gfifteen\) graph.  Although the quotient graph has \(120\)
automorphisms, no relabeling is needed for this comparison: the
lexicographically registered orbit labels already give the historical
labels.

Each four-point orbit contains exactly two orbits of the normal
\(\Ctwo\).  Quotienting first by \(\Ctwo\) and then by
\(\Vfour/\Ctwo\) produces the same thirty \(\Gfifteen\) edges as
quotienting \(\Gsixty\) directly by \(\Vfour\).

\begin{theorem}[Exact quotient factorization]
Conditional on \(\Vfour\) receipt and full-automorphism covariance,
class \(20\) is the unique normal \(\Vfour\) action.  It contains the
normal binary class \(22\), and its orbit maps factor exactly as
\[
  \Gsixty
  \longrightarrow
  \Gsixty/\Ctwo\cong\Gthirty
  \longrightarrow
  \Gsixty/\Vfour\cong\Gfifteen.
\]
The intermediate deck group is
\(\Vfour/\Ctwo\cong\Ctwo\).
\end{theorem}
